% Part: first-order-logic
% Chapter: tableaux
% Section: provability-consistency

\documentclass[../../../include/open-logic-section]{subfiles}

\begin{document}

\iftag{FOL}
      {\olfileid{fol}{tab}{prv}}
      {\olfileid{pl}{tab}{prv}}

\olsection{\usetoken{S}{derivability} आणि सुसंगतता}

आता आपण !!{derivability} संबंधाचे अनेक गुणधर्म सिद्ध करू. ते स्वतंत्रपणेही
महत्त्वाचे आहेत, पण संपूर्णता प्रमेयाच्या सिद्धतेत प्रत्येकाची भूमिका असेल.

\begin{prop}\ollabel{prop:provability-contr}
  जर $\Gamma \Proves !A$ आणि $\Gamma \cup \{!A\}$ विसंगत असेल, तर
  $\Gamma$ विसंगत आहे.
\end{prop}

\begin{proof}
$\Gamma_0 = \{!B_1, \dots, !B_n\}$ आणि $\Gamma_1
  =\{!C_1, \dots, !C_m\} \subseteq \Gamma$ असे सांत उपसंच असतात की
  \begin{align*}
    \{\sFmla{\False}{!A}, &
    \sFmla{\True}{!B_1}, \dots, \sFmla{\True}{!B_n}\} \\
    \{\sFmla{\True}{!A}, &
    \sFmla{\True}{!C_1}, \dots, \sFmla{\True}{!C_m}\}
  \end{align*}
  यांसाठी बंद !!{tableau}s असतात. $!A$ वर \Cut{} नियम वापरून हे दोन्ही
  एकाच बंद !!{tableau} मध्ये एकत्र करता येतात; त्यावरून $\Gamma_0
  \cup \Gamma_1$ विसंगत असल्याचे दिसते. $\Gamma_0
  \subseteq \Gamma$ आणि $\Gamma_1 \subseteq \Gamma$ असल्यामुळे $\Gamma_0 \cup
  \Gamma_1 \subseteq \Gamma$; म्हणून $\Gamma$ विसंगत आहे.
\end{proof}

\begin{prop}
\ollabel{prop:prov-incons}
$\Gamma \Proves !A$ तेव्हा आणि केवळ तेव्हाच, जेव्हा $\Gamma \cup \{\lnot !A\}$
विसंगत असते.
\end{prop}

\begin{proof}
प्रथम $\Gamma \Proves !A$ असे समजा; म्हणजे पुढील संचासाठी बंद
!!{tableau} असतो:
\[
\{\sFmla{\False}{!A},
\sFmla{\True}{!B_1}, \dots, \sFmla{\True}{!B_n}\}
\]
$\TRule{\True}{\lnot}$ नियम वापरून त्याचे पुढील संचासाठीच्या बंद
!!{tableau} मध्ये रूपांतर करता येते:
\[
\{\sFmla{\True}{\lnot !A},
\sFmla{\True}{!B_1}, \dots, \sFmla{\True}{!B_n}\}.
\]

दुसरीकडे, दुसऱ्या संचासाठी बंद !!{tableau} असेल, तर त्यातील पहिले
गृहीतक~$\sFmla{\True}{\lnot !A}$ आणि त्या गृहीतकाला
\TRule{\True}{\lnot} लावून मिळणारी सर्व सूत्रे काढून, तसेच
$\sFmla{\False}{!A}$ हे गृहीतक जोडून, त्याचे पहिल्या संचासाठीच्या बंद
!!{tableau} मध्ये रूपांतर करता येते. कारण एखादी शाखा पूर्वी
$\sFmla{\True}{\lnot !A}$ ला \TRule{\True}{\lnot} लावून मिळणारा
निष्कर्ष, म्हणजे $\sFmla{\False}{!A}$, असल्यामुळे बंद झाली असेल, तर नव्या
!!{tableau} मधील संबंधित शाखाही बंद असते. जुन्या टॅब्लोतील एखादी शाखा
$\sFmla{\True}{\lnot !A}$ हे गृहीतक आणि $\sFmla{\False}{\lnot !A}$
ही दोन्ही सूत्रे असल्यामुळे बंद झाली असेल, तर $\sFmla{\False}{\lnot !A}$
ला $\TRule{\False}{\lnot}$ लावून $\sFmla{\True}{!A}$ मिळवता येते आणि
ती शाखा बंद करता येते. आपण $\sFmla{\False}{!A}$ हे गृहीतक जोडलेले
असल्यामुळे ही शाखा बंद होते.
\end{proof}

\begin{prob}
$\Gamma \Proves \lnot !A$ तेव्हा आणि केवळ तेव्हाच, जेव्हा $\Gamma \cup \{!A\}$
विसंगत असते, हे सिद्ध करा.
\end{prob}

\begin{prop}\ollabel{prop:explicit-inc}
  जर $\Gamma \Proves !A$ आणि $\lnot !A \in \Gamma$ असेल, तर $\Gamma$
  विसंगत आहे.
\end{prop}

\begin{proof}
  समजा $\Gamma \Proves !A$ आणि $\lnot !A \in \Gamma$. मग
  $!B_1$, \dots, $!B_n \in \Gamma$ अशी वाक्ये असतात की
  \[
  \{\sFmla{\False}{!A}, \sFmla{\True}{!B_1}, \dots, \sFmla{\True}{!B_n}\}
  \]
  यासाठी बंद टॅब्लो असतो. \sFmla{\False}{!A} हे गृहीतक
  \sFmla{\True}{\lnot !A} ने बदला आणि \sFmla{\True}{\lnot !A} ला
  \TRule{\True}{\lnot} लावून मिळणारा निष्कर्ष गृहीतकांनंतर घाला. जुन्या
  !!{tableau} मध्ये ओळ~$1$ च्या आधारे समर्थित कोणतेही !!{sentence} आता
  ओळ~$n+2$ च्या आधारे समर्थित होते. म्हणून जुना !!{tableau} बंद असेल, तर
  नवा !!{tableau}सुद्धा बंद असतो. सर्व गृहीतके~$\Gamma$ मध्ये असल्यामुळे हा टॅब्लो
  $\Gamma$ विसंगत असल्याचे दाखवतो.
\end{proof}

\begin{prop}\ollabel{prop:provability-exhaustive}
  $\Gamma \cup \{!A\}$ आणि $\Gamma \cup \{\lnot !A\}$ हे दोन्ही संच
  विसंगत असतील, तर $\Gamma$ विसंगत आहे.
\end{prop}

\begin{proof}
  $!B_1$, \dots, $!B_n \in \Gamma$ आणि $!C_1$, \dots,
  $!C_m \in \Gamma$ अशी वाक्ये असतील आणि
  \begin{align*}
    \{\sFmla{\True}{!A}, &
    \sFmla{\True}{!B_1}, \dots, \sFmla{\True}{!B_n}\} \text{ आणि}\\
    \{\sFmla{\True}{\lnot !A}, &
    \sFmla{\True}{!C_1}, \dots, \sFmla{\True}{!C_m}\}
  \end{align*}
  या दोन्हींसाठी बंद !!{tableau}s असतील, तर
  $\sFmla{\True}{!B_1}$, \dots, $\sFmla{\True}{!B_n}$ आणि
  $\sFmla{\True}{!C_1}$, \dots, $\sFmla{\True}{!C_m}$ ही गृहीतके घेऊन,
  त्यानंतर \Cut{} नियम लावून, $\Gamma$ विसंगत असल्याचे दाखवणारा एकच
  संयुक्त !!{tableau} तयार करता येतो. त्यामुळे दोन शाखा मिळतात: एक
  $\sFmla{\True}{!A}$ ने, तर दुसरी $\sFmla{\False}{!A}$ ने सुरू होते.
  
  डाव्या बाजूस पहिल्या !!{tableau} मधील त्याच्या गृहीतकांखालचा भाग जोडा.
  येथे प्रत्येक नियमाचा उपयोग अजूनही योग्य आहे, कारण
  $\sFmla{\True}{!A}$ सहित पहिल्या !!{tableau} मधील प्रत्येक गृहीतक
  उपलब्ध आहे. त्यामुळे $\sFmla{\True}{!A}$ च्या खालील प्रत्येक शाखा बंद होते.
  
  उजव्या बाजूस दुसऱ्या !!{tableau} मधील त्याच्या गृहीतकांखालचा भाग जोडा;
  मात्र $\sFmla{\True}{\lnot !A}$ ला~$\TRule{\True}{\lnot}$ लावून मिळणारे
  सर्व परिणाम काढून टाका. $\sFmla{\True}{\lnot !A}$ ला
  $\TRule{\True}{\lnot}$ लावल्याचा निष्कर्ष $\sFmla{\False}{!A}$ हा आहे;
  तरीही तो उपलब्ध असतो, कारण तो संयुक्त !!{tableau} च्या उजव्या बाजूवरील
  \Cut{} नियमाचा निष्कर्ष आहे.

  दुसऱ्या टॅब्लोतील एखादी शाखा $\sFmla{\True}{\lnot !A}$ हे गृहीतक
  (जे आता संयुक्त !!{tableau} मध्ये गृहीतक म्हणून येत नाही) आणि
  $\sFmla{\False}{\lnot !A}$ ही दोन्ही सूत्रे असल्यामुळे बंद झाली असेल,
  तर $\sFmla{\False}{\lnot !A}$ ला $\TRule{\False}{\lnot}$ लावून
  $\sFmla{\True}{!A}$ मिळवता येते. मग संयुक्त !!{tableau} मधील संबंधित
  शाखासुद्धा बंद होते, कारण तिच्यात \Cut{} नियमाचा उजवीकडील निष्कर्ष
  $\sFmla{\False}{!A}$ असतो. दुसऱ्या !!{tableau} मधील एखादी शाखा अन्य
  कोणत्याही कारणाने बंद झाली असेल, तर संयुक्त !!{tableau} मधील संबंधित
  शाखासुद्धा बंद होते; कारण जुन्या दुसऱ्या !!{tableau} मधील त्या शाखेवर
  येणारी $\sFmla{\True}{\lnot !A}$ व्यतिरिक्त इतर सर्व !!{signed formula}s
  संयुक्त !!{tableau} मधील संबंधित शाखेवरही येतात.
  \end{proof}

\end{document}
